\begin{summarybox}
\textbf{\textcolor{CSummary}{一句话核心}}
椭圆 $x^2/a^2+y^2/b^2=1$（$a>b>0$）与直线 $y=kx+m$ 相切，当且仅当
\[
 \boxed{m^2=a^2k^2+b^2}.
\]

\textbf{\textcolor{CSummary}{30 秒速记卡}}
\begin{itemize}[leftmargin=*, itemsep=0.35em, topsep=0.25em]
\item \textbf{切线：}$\displaystyle\boxed{y=kx\pm\sqrt{a^2k^2+b^2}}$。
\item \textbf{切点：}$\displaystyle\boxed{P\left(-\frac{a^2k}{m},\frac{b^2}{m}\right)}$（先确认相切）。
\item \textbf{顺带判位置：}比较 $m^2$ 与 $D=a^2k^2+b^2$；小于为割线、等于为切线、大于为相离。
\item \textbf{特殊线：}竖直切线 $x=\pm a$；水平切线 $y=\pm b$。
\end{itemize}

\medskip
\textbf{\textcolor{CSummary}{五种推导，一句话认方法}}
\begin{description}[style=nextline, leftmargin=2.8em, labelsep=0.5em, itemsep=0.35em]
\item[\textbf{最快：伸缩变圆}] $x=au,y=bv$，直线变为 $-aku+bv=m$；圆心距等于 $1$。
\item[\textbf{最通用：判别式}] $\Delta=4a^2b^2(a^2k^2+b^2-m^2)=0$。
\item[\textbf{焦点法：距离积}] $d_1d_2=b^2$，\textbf{同时两焦点须在切线同侧}，得到 $(m^2-c^2k^2)/(1+k^2)=b^2$。
\item[\textbf{切点法：点切式}] $x_0x/a^2+y_0y/b^2=1$，对比 $k,m$ 并代入椭圆方程。
\item[\textbf{最值法：参数式}] $m=b\sin\theta-ak\cos\theta$，相切恰好取到最大值或最小值。
\end{description}

\medskip
\textbf{\textcolor{CSummary}{焦点乘积的关键反例}}
在 $x^2/9+y^2/4=1$ 中，直线 $y=2x$ 经过原点，是割线；但其两焦点到直线的距离均为 $2$，故 $d_1d_2=4=b^2$。\textbf{因此“距离乘积 $=b^2$”单独不是相切的充要条件！}

\medskip
\textbf{\textcolor{CSummary}{和圆的关联}}
圆 $x^2+y^2=R^2$ 的对应式为 $m^2=R^2(k^2+1)$，本质是圆心距等于半径；椭圆切线式就是它经过横、纵方向不同伸缩后的版本。
\end{summarybox}
